% Some basic packages
\usepackage[utf8]{inputenc}
\usepackage[T1]{fontenc}
\usepackage{textcomp}
\usepackage[french]{babel}
\usepackage{url}
\usepackage{graphicx}
\usepackage{float}
\usepackage{booktabs}
\usepackage{enumitem}
\usepackage{enumerate}

\pdfminorversion=7

% Don't indent paragraphs, leave some space between them
\usepackage{parskip}

% Hide page number when page is empty
\usepackage{emptypage}
\usepackage{subcaption}
\usepackage{multicol}
\usepackage{xcolor}

% Other font I sometimes use.
% \usepackage{cmbright}

% Math stuff
\usepackage{amsmath, amsfonts, mathtools, amsthm, amssymb}
% Fancy script capitals
\usepackage{mathrsfs}
\usepackage{cancel}
% Bold math
\usepackage{bm}
% Some shortcuts
\newcommand\A{\ensuremath{\mathbb{A}}}
\newcommand\B{\ensuremath{\mathbb{B}}}
\newcommand\C{\ensuremath{\mathbb{C}}}
\newcommand\D{\ensuremath{\mathbb{D}}}
\newcommand\E{\ensuremath{\mathbb{E}}}
\newcommand\F{\ensuremath{\mathbb{F}}}
\newcommand\G{\ensuremath{\mathbb{G}}}
\newcommand\H{\ensuremath{\mathbb{H}}}
\newcommand\I{\ensuremath{\mathbb{I}}}
\newcommand\J{\ensuremath{\mathbb{J}}}
\newcommand\K{\ensuremath{\mathbb{K}}}
\newcommand\L{\ensuremath{\mathbb{L}}}
\newcommand\M{\ensuremath{\mathbb{M}}}
\newcommand\N{\ensuremath{\mathbb{N}}}
\newcommand\O{\ensuremath{\mathbb{O}}}
\newcommand\P{\ensuremath{\mathbb{P}}}
\newcommand\Q{\ensuremath{\mathbb{Q}}}
\newcommand\R{\ensuremath{\mathbb{R}}}
\newcommand\S{\ensuremath{\mathbb{S}}}
\newcommand\T{\ensuremath{\mathbb{T}}}
\newcommand\U{\ensuremath{\mathbb{U}}}
\newcommand\V{\ensuremath{\mathbb{V}}}
\newcommand\W{\ensuremath{\mathbb{W}}}
\newcommand\X{\ensuremath{\mathbb{X}}}
\newcommand\Y{\ensuremath{\mathbb{Y}}}
\newcommand\Z{\ensuremath{\mathbb{Z}}}

\newcommand\aA{\ensuremath{\mathcal{A}}}
\newcommand\bB{\ensuremath{\mathcal{B}}}
\newcommand\cC{\ensuremath{\mathcal{C}}}
\newcommand\dD{\ensuremath{\mathcal{D}}}
\newcommand\eE{\ensuremath{\mathcal{E}}}
\newcommand\fF{\ensuremath{\mathcal{F}}}
\newcommand\gG{\ensuremath{\mathcal{G}}}
\newcommand\hH{\ensuremath{\mathcal{H}}}
\newcommand\iI{\ensuremath{\mathcal{I}}}
\newcommand\jJ{\ensuremath{\mathcal{J}}}
\newcommand\kK{\ensuremath{\mathcal{K}}}
\newcommand\lL{\ensuremath{\mathcal{L}}}
\newcommand\mM{\ensuremath{\mathcal{M}}}
\newcommand\nN{\ensuremath{\mathcal{N}}}
\newcommand\oO{\ensuremath{\mathcal{O}}}
\newcommand\pP{\ensuremath{\mathcal{P}}}
\newcommand\qQ{\ensuremath{\mathcal{Q}}}
\newcommand\rR{\ensuremath{\mathcal{R}}}
\newcommand\sS{\ensuremath{\mathcal{S}}}
\newcommand\tT{\ensuremath{\mathcal{T}}}
\newcommand\uU{\ensuremath{\mathcal{U}}}
\newcommand\vV{\ensuremath{\mathcal{V}}}
\newcommand\wW{\ensuremath{\mathcal{W}}}
\newcommand\xX{\ensuremath{\mathcal{X}}}
\newcommand\yY{\ensuremath{\mathcal{Y}}}
\newcommand\zZ{\ensuremath{\mathcal{Z}}}

\newcommand{\im}{\operatorname{im}}
\newcommand{\rang}{\operatorname{rang}}
\newcommand{\End}{\operatorname{End}}
\newcommand{\id}{\operatorname{id}}
\newcommand{\GL}{\operatorname{GL}}
\newcommand{\Bil}{\operatorname{Bil}}
\newcommand{\Bij}{\operatorname{Bij}}
\newcommand{\DET}{\operatorname{det}}
\newcommand{\min}{\operatorname{min}}
\newcommand{\max}{\operatorname{max}}
\newcommand{\sup}{\operatorname{sup}}
\newcommand{\TR}{\operatorname{Tr}}
\newcommand{\inf}{\operatorname{inf}}
\newcommand{\cof}{\operatorname{Cof}}
\renewcommand{\span}{\operatorname{Vect}}
\renewcommand{\inte}{\operatorname{int}}
\newcommand{\const}{\textsc{const}}

\renewcommand{\d}{\text{d}}
\let\p\partial


% Easily typeset systems of equations (French package)
\usepackage{systeme}

% Put x \to \infty below \lim
\let\svlim\lim\def\lim{\svlim\limits}
% Idem for the sum
\let\svsum\sum\def\sum{\svsum\limits}
% Idem for the min
\let\svmin\min\def\min{\svmin\limits}
% Idem for the max
\let\svmax\max\def\max{\svmax\limits}
% And the integral
\let\svint\int\def\int{\svint\limits}
% And the fractions
\let\svfrac\frac\def\frac#1#2{{\displaystyle\svfrac{#1}{#2}}}


%Make implies and impliedby shorter
\let\implies\Longrightarrow
\let\impliedby\Longleftarrow
\let\iff\Longleftrightarrow
\let\bf\boldsymbol
\let\l\left
\let\r\right

% Add \contra symbol to denote contradiction
\usepackage{stmaryrd} % for \lightning
\newcommand\contra{\scalebox{1.5}{$\lightning$}}

\let\phi\varphi
\let\epsilon\varepsilon

% Command for short corrections
% Usage: 1+1=\correct{3}{2}

\definecolor{correct}{HTML}{009900}
\newcommand\correct[2]{\ensuremath{\:}{\color{red}{#1}}\ensuremath{\to }{\color{correct}{#2}}\ensuremath{\:}}
\newcommand\green[1]{{\color{correct}{#1}}}

% horizontal rule
\newcommand\hr{
    \noindent\rule[0.5ex]{\linewidth}{0.5pt}
}

% hide parts
\newcommand\hide[1]{}

% si unitx
\usepackage{siunitx}
\sisetup{locale = FR}

% Environments
\makeatother
% For box around Definition, Theorem, \ldots
\usepackage{mdframed}
\mdfsetup{skipabove=1em,skipbelow=0em}
\theoremstyle{definition}
%\newmdtheoremenv[nobreak=true]{definition}{Définition}
%\newmdtheoremenv[nobreak=true]{propriete}{Propriété}
%\newmdtheoremenv[nobreak=true]{corollaire}{Corollaire}
%\newmdtheoremenv[nobreak=true]{lemme}{Lemme}
%\newmdtheoremenv[nobreak=true]{proposition}{Proposition}
%\newmdtheoremenv[nobreak=true]{theoreme}{Théorème}
%\newmdtheoremenv[nobreak=true]{loi}{Loi}
%\newmdtheoremenv[nobreak=true]{postulat}{Postulat}
\newmdtheoremenv{conclusion}{Conclusion}
\newmdtheoremenv{bonus}{Supplément}
\newmdtheoremenv{conjecture}{Conjecture}

\newtheorem*{rappel}{Rappel}
%\newtheorem*{preuve}{Preuve}
\newtheorem*{notation}{Notation}
\newtheorem*{observation}{Observation}
\newtheorem*{reflexion}{Réflexion}
\newtheorem*{exercice}{Exercice}
\newtheorem*{remarque}{Remarque}
\newtheorem*{pratique}{En pratique}
\newtheorem*{probleme}{Problème}
\newtheorem*{terminologie}{Terminologie}
\newtheorem*{application}{Application}
\newtheorem*{exemple}{Exemple}
\newtheorem*{question}{Question}
\newtheorem*{zoologie}{Zoologie}
\newtheorem*{idee}{Idée}

% Personal environement
\newtheoremstyle{bleuPreuve}     % nom du style
  {}{}                           % espace avant/après
  {\color{blue!60!black}}        % corps du texte : légèrement bleuté
  {}                             % indentation
  {\itshape\color{blue!70!black}}% tête du théorème ("Preuve") : italique et bleutée
  {.}                            % ponctuation après le titre
  { }                            % espace après le titre
  {}                             % définition du titre

\theoremstyle{bleuPreuve}
\newtheorem*{preuve}{Preuve}

\newtheoremstyle{vertRetour}     % nom du style
  {}{}                           % espace avant/après
  {\color{green!40!black}}        % corps du texte : légèrement bleuté
  {}                             % indentation
  {\itshape\color{black}}% tête du théorème ("Preuve") : italique et bleutée
  {\\}                            % ponctuation après le titre
  { }                            % espace après le titre
  {}                             % définition du titre

\theoremstyle{vertRetour}
\newtheorem*{retour}{}

% End example and intermezzo environments with a small diamond (just like proof
% environments end with a small square)
\usepackage{etoolbox}
\AtEndEnvironment{preuve}{\null\hfill$\square$}%
%\AtEndEnvironment{intermezzo}{\null\hfill$\square$}%
% \AtEndEnvironment{opmerking}{\null\hfill$\diamond$}%

% Fix some spacing
% http://tex.stackexchange.com/questions/22119/how-can-i-change-the-spacing-before-theorems-with-amsthm
\makeatletter
\def\thm@space@setup{%
  \thm@preskip=\parskip \thm@postskip=0pt
}


% Exercise 
% Usage:
% \oefening{5}
% \suboefening{1}
% \suboefening{2}
% \suboefening{3}
% gives
% Oefening 5
%   Oefening 5.1
%   Oefening 5.2
%   Oefening 5.3
\newcommand{\oefening}[1]{%
    \def\@oefening{#1}%
    \subsection*{Oefening #1}
}

\newcommand{\suboefening}[1]{%
    \subsubsection*{Oefening \@oefening.#1}
}


% \lecture starts a new lecture (les in dutch)
%
% Usage:
% \lecture{1}{di 12 feb 2019 16:00}{Inleiding}
%
% This adds a section heading with the number / title of the lecture and a
% margin paragraph with the date.

% I use \dateparts here to hide the year (2019). This way, I can easily parse
% the date of each lecture unambiguously while still having a human-friendly
% short format printed to the pdf.

\usepackage{xifthen}
\def\testdateparts#1{\dateparts#1\relax}
\def\dateparts#1 #2 #3 #4 #5\relax{
    \marginpar{\small\textsf{\mbox{#1 #2 #3 #5}}}
}

\def\@lecture{}%
\newcommand{\lecture}[3]{
    \ifthenelse{\isempty{#3}}{%
        \def\@lecture{Session #1}%
    }{%
        \def\@lecture{Session #1: #3}%
    }%
    \subsection*{\@lecture}
    \marginpar{\small\textsf{\mbox{#2}}}
}



% These are the fancy headers
\usepackage{fancyhdr}
\pagestyle{fancy}

% LE: left even
% RO: right odd
% CE, CO: center even, center odd
% My name for when I print my lecture notes to use for an open book exam.
% \fancyhead[LE,RO]{Gilles Castel}

\fancyhead[RO,LE]{\@lecture} % Right odd,  Left even
\fancyhead[RE,LO]{}          % Right even, Left odd

\fancyfoot[RO,LE]{\thepage}  % Right odd,  Left even
\fancyfoot[RE,LO]{}          % Right even, Left odd
\fancyfoot[C]{\leftmark}     % Center

\makeatother




% Todonotes and inline notes in fancy boxes
\usepackage{todonotes}
\usepackage[most]{tcolorbox}

% Make boxes breakable
\tcbuselibrary{skins, breakable}

% Verbetering is correction in Dutch
% Usage: 
% \begin{verbetering}
%     Lorem ipsum dolor sit amet, consetetur sadipscing elitr, sed diam nonumy eirmod
%     tempor invidunt ut labore et dolore magna aliquyam erat, sed diam voluptua. At
%     vero eos et accusam et justo duo dolores et ea rebum. Stet clita kasd gubergren,
%     no sea takimata sanctus est Lorem ipsum dolor sit amet.
% \end{verbetering}

% === Style général ===
\tcbset{
  definitionstyle/.style={
    enhanced,
    colback=white, 
    breakable, 
    colframe=green!60!black,
    colbacktitle=green!60!black, 
    coltitle=white,
    fonttitle=\bfseries\sffamily,
    boxrule=0.8pt, 
    arc=0mm,
    before skip=1em, 
    after skip=1em,
    title=#1
  },
  theoremstyle/.style={
    enhanced,
    colback=white, 
    breakable, 
    colframe=blue!75!black,
    colbacktitle=blue!75!black, 
    coltitle=white,
    fonttitle=\bfseries\sffamily,
    boxrule=0.8pt, 
    arc=0mm,
    before skip=1em, 
    after skip=1em,
    title=Théorème~#1
  },
  lemmestyle/.style={
    enhanced,
    colback=white, 
    breakable, 
    colframe=blue!50!black,
    colbacktitle=blue!50!black, 
    coltitle=white,
    fonttitle=\bfseries\sffamily,
    boxrule=0.8pt, 
    arc=0mm,
    before skip=1em, 
    after skip=1em,
    title=Lemme~#1
  },
  corollairestyle/.style={
    enhanced,
    colback=white, 
    breakable, 
    colframe=blue!100!black,
    colbacktitle=blue!100!black, 
    coltitle=white,
    fonttitle=\bfseries\sffamily,
    boxrule=0.8pt, 
    arc=0mm,
    before skip=1em, 
    after skip=1em,
    title=Corollaire~#1
  },
  propositionstyle/.style={
    enhanced,
    colback=white,
    breakable, 
    colframe=blue!70!red,
    colbacktitle=blue!70!red, 
    coltitle=white,
    fonttitle=\bfseries\sffamily,
    boxrule=0.8pt, 
    arc=0mm,
    before skip=1em, 
    after skip=1em,
    title=Proposition~#1
  },
  proprietestyle/.style={
    enhanced,
    colback=white, 
    breakable, 
    colframe=blue!60!red,
    colbacktitle=blue!60!red, 
    coltitle=white,
    fonttitle=\bfseries\sffamily,
    boxrule=0.8pt, 
    arc=0mm,
    before skip=1em, 
    after skip=1em,
    title=Propriété~#1
  },
  loistyle/.style={
    enhanced,
    colback=white, 
    breakable, 
    colframe=yellow!60!red,
    colbacktitle=yellow!60!red, 
    coltitle=white,
    fonttitle=\bfseries\sffamily,
    boxrule=0.8pt, 
    arc=0mm,
    before skip=1em, 
    after skip=1em,
    title=Loi~#1
  },
  postulatstyle/.style={
    enhanced,
    colback=white, 
    breakable, 
    colframe=yellow!80!red,
    colbacktitle=yellow!80!red, 
    coltitle=white,
    fonttitle=\bfseries\sffamily,
    boxrule=0.8pt, 
    arc=0mm,
    before skip=1em, 
    after skip=1em,
    title=Postulat~#1
  }
}


% === Environnement principal ===
\newenvironment{definition}[1][]
{\begin{tcolorbox}[definitionstyle={#1}]}
{\end{tcolorbox}}

\newenvironment{theoreme}[1][]
{\begin{tcolorbox}[theoremstyle={#1}]}
{\end{tcolorbox}}

\newenvironment{lemme}[1][]
{\begin{tcolorbox}[lemmestyle={#1}]}
{\end{tcolorbox}}

\newenvironment{corollaire}[1][]
{\begin{tcolorbox}[corollairestyle={#1}]}
{\end{tcolorbox}}

\newenvironment{proposition}[1][]
{\begin{tcolorbox}[propositionstyle={#1}]}
{\end{tcolorbox}}

\newenvironment{propriete}[1][]
{\begin{tcolorbox}[proprietestyle={#1}]}
{\end{tcolorbox}}

\newenvironment{loi}[1][]
{\begin{tcolorbox}[loistyle={#1}]}
{\end{tcolorbox}}

\newenvironment{postulat}[1][]
{\begin{tcolorbox}[postulatstyle={#1}]}
{\end{tcolorbox}}

% === Sous-environnement : Contexte Hypothèse main ===
\newenvironment{contexte}
{
  \par\vspace{0.5em}\noindent\textit{Soit:} \par\vspace{0.3em}
  \begin{tcolorbox}[
    enhanced, 
    colback=black!5, 
    colframe=black!5,
    boxrule=0.4pt, 
    arc=0.5mm, 
    left=1mm, 
    right=1mm,
    before skip=2pt, 
    after skip=4pt, 
    sharp corners,
    breakable]
}
{
  \end{tcolorbox}
}

\newenvironment{hypothese}
{
  \par\vspace{0.5em}\noindent\textit{Supposons:} \par\vspace{0.3em}
  \begin{tcolorbox}[
    enhanced, 
    colback=black!5, 
    colframe=black!5,
    boxrule=0.4pt, 
    arc=0.5mm, 
    left=1mm, 
    right=1mm,
    before skip=2pt, 
    after skip=4pt, 
    sharp corners,
    breakable]
}
{
  \end{tcolorbox}
}

\newenvironment{main}
{\par\vspace{0.5em}}
{\par\vspace{0.5em}}

\newenvironment{donnee}
{\itshape}
{}

\newenvironment{meta}{
  \par\vspace{0.5em}
  \begin{tcolorbox}[
    enhanced,
    colback=white!95!black,             % fond blanc
    colframe=gray!70,          % couleur de la bordure (barre)
    left=8mm, 
    colframe=white,                  % espace entre la barre et le texte
    right=0mm, top=0mm, bottom=0mm,
    boxrule=0.8pt,             % épaisseur de la barre
    borderline west={2pt}{0pt}{gray!70}, % barre verticale à gauche
    sharp corners,              % coins non arrondis
    boxsep=0pt,
    before skip=6pt, after skip=6pt,
    halign=center              % texte centré
  ]
  \itshape \vspace{0.5em}                    % texte en italique
}{
  \vspace{0.5em} \end{tcolorbox}\par\vspace{0.5em}
}

% sigure support as explained in my blog post.
\usepackage{import}
\usepackage{xifthen}
\usepackage{pdfpages}
\usepackage{transparent}
\newcommand{\incfig}[1]{%
    \def\svgwidth{\columnwidth}
    \import{./figures/}{#1.pdf_tex}
}

% Fix some stuff
% %http://tex.stackexchange.com/questions/76273/multiple-pdfs-with-page-group-included-in-a-single-page-warning
\pdfsuppresswarningpagegroup=1


% My name
\author{Emiliem Progin}
